\subsection{The pace, and where the cost went}

An account of this method that did not state its speed would be evasive, since
speed is the whole of what changed. The figures are in the repository's own
history and can be recovered from it.

The transcription of 667 sheets took nine days, from 5 to 13 September 2026, in
sixty-two passes of up to twelve sheets. Nobody keyed a character: a pass is a
photograph read and a file written, and the twelve sheets of a batch take the
model minutes. In the same fortnight the repository received the scripts that derive the
accounts, the formats, the materials and the calendar (some 11,000 lines), the
site that places the reading beside the photograph (some 12,500 lines), the
TEI converter and its validator, and later the 1,800 lines of Lean that
certify the arithmetic. A transcription project of this size, keyed and
proofed by people, is measured in a person's months; this was a fortnight of
one person's attention and a model's time.

That comparison is the easy half, and it is not the interesting one. What the
method actually changes is the second operation, the one that has no name in
the traditional workflow: \emph{restructuring}. Once the transcription is a
constrained language rather than a document, every question one might put to
the archive becomes a script of a few hundred lines --- what was paid and by
whom, which formats recur, when the colourman changed, on what days the books
were written in. Each of those is a census this edition publishes, each was
written in an afternoon, and each is regenerated on every build. The cost of
asking a new question of the corpus has fallen by more than the cost of reading
it, because reading is done once and asking is done for ever.

The cost did not disappear; it moved, and it moved to the two places this
article has been describing. The first is the definition of rules. The Book~IV
reader changed its rule four times across the volume's fifty-four years, and
every one of those changes is an editorial decision about what the writer meant
that no model and no script could take: whether a bare figure is a subtotal,
whether a \q{less} line is a deduction, whether three payments are three sales
or one. Writing them down --- in the method note, then as Lean functions proved
for every input --- took longer than transcribing the volume. The second is
verification. 667 sheets are read and none is checked. A pass costs minutes and
a human check costs an hour a batch, so the reading has produced, at the pace
described above, a debt of some fifty hours that the method cannot discharge
and does not pretend to.

The same arithmetic applies to the three sibling editions of Section~8, and
states the size of what is left. The Grothendieck fonds stands at 1,937 pages
in 122 passes, sixteen pages to a pass, against the sixteen thousand pages the
edition's own catalogue counts as openly accessible at Montpellier: at that
rate the reading is some nine hundred further passes, which is months of
machine time and, at an hour a batch, about nine hundred hours of human
checking --- that is, a reading a machine can finish this year and a
verification no individual can. The mathematicians' manuscripts in Gallica
stand at 647 views in 35 passes, eighteen views to a pass, against 25,258
views catalogued: some 1,400 further passes. Neither number is an argument for doing it; both are the
argument of this article, which is that the transcription is no longer the
bottleneck and the checking is, and that an edition which cannot say which of
its sheets a person has read is publishing a machine's confidence as though it
were scholarship.

There is a third cost the traditional account treats as fixed, and it is the
apparatus itself. Everything in Figure~\ref{fig:pipeline} below the reading
--- the converters, the censuses, the validators, the site, the workflows
that rebuild and compare --- is code, and code of that kind was the part of
an edition that took a developer and a grant. Here it was written by the same
model that read the sheets, through a code assistant, under one person's
direction, at a pace the repositories' histories give. This edition's
apparatus, some twenty-three thousand lines of scripts and interface and
eighteen hundred of Lean, was written between 5 and 20 September 2026 in some two
hundred and twenty commits. The siblings show what an edition costs once the schema exists: the
Grothendieck edition, first and therefore longest, 14,500 lines in 246
commits over six weeks; the Gallica edition, 11,800 lines in 56 commits over
ten days; the morphology course, 5,700 lines in 104 commits over two days,
the schema by then known and only its values new. Lines and commits measure
activity, not quality; what they establish is that a stated schema can be
filled in for a new corpus in days.

That ease has two consequences, both met here. Code written at that pace is
written faster than it is read, and what makes it trustworthy is the
discipline described for the transcription: every derived file regenerated
and compared on every build, every rule stated where a reader can find it, a
build that fails loudly. Two of the reader's rules were found by bugs
(Section 6.2), and the calendar of Figure~\ref{fig:days} went stale because its
generator wrote to a path no build checked, until the comparison caught it. The
assistant does not remove the checks; it makes them the whole of the
guarantee, since nobody reads the code as an editor reads a leaf. And when a
converter or a census costs an afternoon, what remains expensive are the
editorial decisions --- what a bare figure is, whether three payments are
one sale, where a colour convention begins --- which the code cannot take.
The assistant moves the cost of an edition further toward its rules and its
checking, which is where it belongs.
